\documentclass[class=scrbook, crop=false]{standalone}
\usepackage[subpreambles=true]{standalone}
\ifstandalone
    % WARNING: Proceed with caution!

% -----------------------------------------------------------------------------------
% For package standalone
% -----------------------------------------------------------------------------------
\usepackage{import}

% -----------------------------------------------------------------------------------
% Language and typeset
% -----------------------------------------------------------------------------------
\usepackage[ngerman, english]{babel}

% Umlauts and other special characters (UTF-8)
% \usepackage[utf8]{inputenc}
\usepackage{fontspec}
\setsansfont{Arial}
% \usepackage[T1]{fontenc}  % Enable accented characters and umlauts
% LuaLatex doesn't need fontenc and uses UTF-8
% \usepackage{lmodern}  % Font face


% --------------------------------------------------------------------------------
% Page formatting
% --------------------------------------------------------------------------------
% Change the header/footer for chapter beginnings and normal pages
\usepackage[automark,headsepline]{scrlayer-scrpage}

% The package provides an easy and flexible user interface to customize the page
% layout, implementing auto-centering and auto-balancing mechanisms
% WARNING: WHEN CHANGING BCOR (Binding correction), the cover needs reworking!...
\newcommand{\theBCOR}{15mm}  % Define binding correction
\usepackage[
    bindingoffset=\theBCOR,
    % showframe, % Show boxes which indicate margins and paddings
    bottom = 3.5cm, % Margins
      left = 2.5cm,
     right = 2.5cm
] {geometry}

% The package 'float' provides a container for document objects which can not be
% broken over pages, such as tables and figures
% Needed for table and figure indexes  
\usepackage{float}

% support for landscape layout
\usepackage{lscape}

% support of \tablenotes command to add notes under table
\usepackage{threeparttable}

% To allow drawing more professional tables
\usepackage{booktabs}

% --------------------------------------------------------------------------------
% Contents
% --------------------------------------------------------------------------------
% Vector graphics (for Cover page)
\usepackage{tikz} 

% Allows additional parameters when including images
\usepackage{graphicx}

% Roman font family for all headings
\addtokomafont{disposition}{\rmfamily}

% Set the line spacing to 1.5
\usepackage[onehalfspacing]{setspace}

% Improves overall text spacing
% http://www.khirevich.com/latex/microtype/
\usepackage[stretch=10]{microtype}

% Math symbols like mu outside the math environment
\usepackage{textcomp}

% A comprehensive (SI) units package∗
% For defining SI units
\usepackage[
    range-units=single,         % Formatting ranges with single unit indication: 1 - 2 m
    range-phrase=-,             % Phrase for range: 1 - 2 m vs 1 to 2 m
    separate-uncertainty=true,  % sets +- between value and uncertainty 
    multi-part-units=repeat     % In expressions with multiple values (multi part numbers) 
                                % the unit is printed each time: 1 mm x 1 mm
] {siunitx}
% https://tex.stackexchange.com/questions/124488/multi-part-numbers-and-units-in-siunitx

% Allows Sourcecodes with highlighting 
\usepackage{listings}

% This package provides user control over the layout of the three basic list
% environments: enumerate, itemize and description
\usepackage{enumitem}
\setlist{nosep} % Remove the vertical space between \item elements in all lists

% ToDo Notes
% \setlength{\marginparwidth}{2cm}
\usepackage{todonotes}
\setuptodonotes{inline, inlinepar}
\reversemarginpar  % Put ToDo notes on the binding's side
% \usepackage{soul} % Colorful ToDo notes

% Check out colors here http://latexcolor.com/
\usepackage{xcolor}

\usepackage{amsmath}    % alignment of equations

% --------------------------------------------------------------------------------
% Other elements
% --------------------------------------------------------------------------------
% Blindtext: Organic looking text dummy
\usepackage{blindtext}

% Hyperlinks within the document (PDF)
% "hidelinks" hides visual highlighting of links
\usepackage[hidelinks]{hyperref}

% Package for Glossary and Index (Acronyms are listed in a separate list) 
\usepackage[acronym, nogroupskip]{glossaries}[=v4.49] % groupskip: alphabetic grouping of entries

\usepackage{xltabular}   % <------- FOR glossaries

% Integration and management of bibliographies
\usepackage{csquotes}   % backend=biber in biblatex needs this package
\usepackage[
    style=ieee,   % style of the bibliography, entries are sorted in alphabetic order. "ieee" is another common style.
    backend=biber,      % based on package 'biber' 
    bibencoding=utf8   % ASCII Text encoding; may use "utf8" instead
] {biblatex}

% --------------------------------------------------------------------------------
%                               PATHS & FILES
% --------------------------------------------------------------------------------
% Fix paths for standalone compiling
%\def \home {D://MS/Thesis/master-thesis}

% Package: scrlayer-scrpage
% \def \stylePath {\home/settings+/style/page}
\input{D://MS/Thesis/master-thesis/settings+/style/page}  % Load page style

% Package: graphicx
\graphicspath{{D://MS/Thesis/master-thesis/images/}}  % Set path to images

% Package: listings
\input{D://MS/Thesis/master-thesis/settings+/style/code.tex}  % Set path to style file
\lstset{inputpath={D://MS/Thesis/master-thesis/code/}} % Default path to code listings

% Package: glossaries
\input{D://MS/Thesis/master-thesis/settings+/style/symbols}  % Set path to symbols list style file
\input{D://MS/Thesis/master-thesis/settings+/style/acronyms}  % Set path to acronym list style file
\input{D://MS/Thesis/master-thesis/references/glossary_acronyms}  % Load glossary, symbol and acronyms list

% Package: biblatex
\addbibresource{D://MS/Thesis/master-thesis/references/references.bib} % Set path to bib resources

% Custom variables
\input{D://MS/Thesis/master-thesis/settings+/variables}
% --------------------------------------------------------------------------------
%                                   OPTIONAL
% --------------------------------------------------------------------------------


% Simple arithmetic for LaTeX commands
% \usepackage{calc}

% Document Elements
% -------------------

% Index
% \usepackage{imakeidx}

% compact Lists
%\usepackage{paralist}

% visual improvements for citations
% \usepackage{epigraph}

% Create pseudo code
% https://www.overleaf.com/learn/latex/Algorithms
% \usepackage{algorithm}
% \usepackage{algorithmic}
%\usepackage[noend]{algpseudocode}

% Formatting
% -------------------
% Tweaks for scrbook, redefines commands of other packages
% \usepackage{scrhack}

% Intelligent space separator (nice for superscript?)
% \usepackage{xspace}

% Allows breaks within tables
%\usepackage{tabularx}

% Allows for page breaks in tables
% \usepackage{longtable}

% allows modifying of captions
% \usepackage{caption}

% Multiline comments
%\usepackage{verbatim}

% % Custom colors
% \definecolor{dartmouthgreen}{rgb}{0.05, 0.5, 0.06}

% IF you want to define unicode characters
% \DeclareUnicodeCharacter{0229}{\c{e}}
% \DeclareUnicodeCharacter{0306}{\u{Z}}


% Document elements
% ------------------------------------

% Table package
% \usepackage{booktabs}

% Pie diagram
% \usepackage{datapie}

% Side by Side images
% \usepackage{subcaption}

% For landscape tables
%\usepackage{pdflscape}
%\usepackage{afterpage}

% Graphics can be flow around by text
%\usepackage{wrapfig}

\fi

% ----------------------------------------------------------------------------
%                                Introduction
% ----------------------------------------------------------------------------
\begin{document}

\chapter{Introduction}
\label{Chapter::Introduction}
    %This chapter outlines the motivation for this thesis

\section{Motivation}
\label{Section::Motivation}
   
     Creep is a time-dependent deformation of materials under sustained stress at elevated temperatures. It plays a critical role in industries such as aerospace(eg: Turbine blades and vanes in jet engine), energy(eg: high temperature pipelines), and manufacturing(eg: concrete structures). Understanding and mitigating creep behavior is essential for ensuring the long-term reliability of components subjected to high-temperature environments. 
     
     Prediction of creep behavior and strength of materials at high temperatures is crucial in the design of structures used at elevated temperatures. For instance, oil vessels in a refinery have to be resistant against the stresses caused by hot fluids for long durations.[M.E. Kassner, M.T. Pérez-Prado
     Erratum to “five-power-law creep in single phase metals and alloys” [progr. mater. sci. 45 (2000) 1–102]
     Prog. Mater. Sci., 45 (2000), p. 273]
     Traditional methods for analyzing creep behavior faces significant challenges as experimental creep tests are time consuming (few months to years). It incurs substantial costs and often yields limited data. Furthermore, numerical simulations using traditional methods can be computationally intensive, demanding fine meshes and iterative solvers that can take hours or even days to run, hindering timely and comprehensive analysis of creep in engineering components. These limitations underscore the need for more efficient and predictive approaches to address the critical role of creep in various industries.
    
    Recent advances in \textit{machine learning (ML)} and \textit{numerical simulation} provide a transformative opportunity to reduce the experimental burden while increasing predictive power. ML models, trained on either simulated or experimental datasets, can act as \textit{surrogate models} that rapidly approximate complex physical processes such as creep. However, building these models effectively requires a thoughtful design of experiments, robust data generation and preprocessing, and intelligent sampling methods, such as \textit{active learning}. Data-driven models are increasingly used due to their straightforward methodology, precise predictive capabilities, reduced costs, and high scalability across different materials and conditions.[ Nava et al. (2023)]
    
    This thesis is motivated by the need to combine \textit{physics-based numerical models} with \textit{data-driven machine learning} to accurately and efficiently predict the creep behavior of metallic materials.    
    
\section{Research Context \& AID4GREENEST Project}
\label{Section::Research_Context_&_AID4GREENEST_Project}
    %\Blindtext
    This work is carried out in the context of the \textbf{AID4GREENEST} project, which is funded by the European Union (EU). It aims to developrange of new Artificial Intelligence-based rapid characterisation methods and modeling tools for the steel sector.["https://aid4greenest.eu/about/"] A core goal of the project is to reduce the cost and time of materials qualification by integrating advanced simulations, smart data sampling, and machine learning.
    
    Specifically, the \texttt{Lagamine} finite element software, developed at the University of Liège, is used to simulate the creep behavior of metallic alloys under varying loads, temperatures, and microstructural parameters. The generated data forms the backbone of this study’s ML training dataset. By combining domain knowledge from materials science with data science techniques, AID4GREENEST seeks to support faster and more sustainable innovation in material development.\\
    
     
    \noindent This thesis serves as a prestudy by:
    
    \begin{itemize}
	    \item Exploring efficient representations of creep curves from standard tests.
	    \item Developing data-driven techniques to create an optimal training dataset.
	    \item Implementing ML-based surrogate models to replicate the simulation outputs efficiently.
	\end{itemize}
	
\section{Research Objectives}
\label{Section::Research_Objective}
	
	The key objectives of this thesis include:
	
	\begin{itemize}
		\item Conducting a literature review on creep behavior, ML applications, and numerical simulation techniques.
		\item Preprocessing and exploring the design space of creep test data.
		\item Investigating different mathematical representations for creep test curves.
		\item Developing a classification model to define valid simulation parameter ranges.
		\item Implementing active learning strategies to optimize data selection.
		\item Training and evaluating ML regression models to predict creep curves.
	\end{itemize}
	
\subsection{Structure of the Thesis}
\label{Section::Structure_of_the_Thesis}

The remainder of this thesis is structured as follows:

\begin{itemize}
	\item \textbf{Chapter 2:} \\
	Discusses the theory of creep in metals, previous modeling approaches, and recent developments in machine learning for material science.
	
	\item \textbf{Chapter 3:} \\
	Explains the numerical simulation setup, data generation, preprocessing steps, sampling strategies, and ML models used.
	
	\item \textbf{Chapter 4:} \\
	Presents model performance, comparisons of sampling and feature extraction strategies, and insights on surrogate modeling.
	
	\item \textbf{Chapter 5:} \\
	Summarizes key contributions and outlines directions for extending this research to real materials and hybrid ML-physics models.
\end{itemize}

\ifstandalone
    \printglossary
    \printbibliography[heading=bibintoc]
\fi

\end{document}
